\documentclass[12pt]{article}

\usepackage[utf8]{inputenc}
\usepackage[spanish]{babel}
\usepackage[shortlabels]{enumitem}

\usepackage{mathtools}
\usepackage{amssymb}
\usepackage{amsthm}

\usepackage{geometry}
\geometry{margin=2.5cm}

\usepackage{hyperref}

\begin{document}

\section*{Ejercicio 1}
Calcula las siguientes derivadas usando la regla de la cadena:
\begin{enumerate}[label=\alph*)]
    \item \(f(x)=\sin(\ln(4x+5))\).
    \item \(g(x)=\arctan\!\big((4x^2+7x+1)^6\big)\).
\end{enumerate}

\section*{Ejercicio 2}
Obtén la ecuación de la recta tangente a
\[
f(x)=\ln\big((x^2-3)^4\big)
\]
en \(x=2\).

\section*{Ejercicio 3}
Estudia la presencia de asíntotas, el crecimiento y la curvatura de la función
\[
f(x)=x+\frac{1}{x}.
\]
Haz un esbozo aproximado de su gráfica.

\section*{Ejercicio 4}
Estudia la derivabilidad de la función
\[
f(x)=
\begin{cases}
x^2\sin\!\left(\frac{1}{x}\right), & x\neq 0,\\[4pt]
0, & x=0,
\end{cases}
\]
en \(x=0\).

\section*{Ejercicio 5}
Calcula la quinta derivada de
\[
f(x)=\sin(x)\cos(x).
\]

\section*{Ejercicio 6}
Hallar \(a\), \(b\) y \(c\) para que la función
\[
f(x)=ax^2+bx+c
\]
pase por el punto \((1,3)\) y sea tangente a la recta
\[
x-y+1=0
\]
en el punto \((2,3)\).

\section*{Ejercicio 7}
Un coche circula por la noche por un camino dado por la curva
\[
y=-x^2+4x+5
\]
de tal forma que el sentido de circulación viene definido por los valores crecientes de la variable \(x\). ¿En qué punto del camino iluminará con sus faros el objeto situado en el punto \((5,4)\)?

Repite el problema suponiendo que el sentido de circulación está definido por los valores decrecientes de la variable \(x\).

\section*{Ejercicio 8}
Calcula la derivada de las siguientes funciones:
\begin{enumerate}[label=\alph*)]
    \item \(f(x)=x^x\).
    \item \(g(x)=(x^2+3)^{6x}\).
\end{enumerate}

\section*{Ejercicio 9}
Consideramos una caja de cartón sin tapa superior y de base cuadrada. Si la suma del área de todas las caras de la caja es \(c^2\) (con \(c>0\)), calcula el volumen máximo de la caja.

\section*{Ejercicio 10}
Utiliza el teorema de Rolle para demostrar que la ecuación
\[
2x^5+x+5=0
\]
no puede tener más de una solución real.

\section*{Ejercicio 11}
Sin calcular la derivada de la función
\[
f(x)=(x-1)(x-2)(x-3)(x-4),
\]
determina cuántas raíces tiene \(f'(x)\) e indica en qué intervalos están.

\section*{Ejercicio 12}
Calcula la ecuación de la recta tangente a
\[
f(x)=\sin\left(\frac{x}{2}\right)
\]
para
\[
x=\frac{2\pi}{3}.
\]

\section*{Ejercicio 13}
Dada la función
\[
f(x)=ax^2+2x+1,
\]
hallar el valor de \(a\) para que la recta tangente a \(f(x)\) en \(x=-1\) sea horizontal.

\section*{Ejercicio 14}
Estudia la continuidad y la derivabilidad de la función
\[
f(x)=
\begin{cases}
-x^2+2, & x<1,\\[4pt]
-2\ln(x)+1, & x\ge 1,
\end{cases}
\]
en \(x=1\).

\end{document}